\chapter{El terremoto: 2025--2026}\label{ch:2026}

\begin{roadmap}
El capítulo~\ref{ch:estructura} cerró con un muro: recuperar la clave
privada costaba, en 2022, unos $2^{9116{,}8}$ frente a los
$2^{140{,}8}$ de la decodificación genérica. Este capítulo cuenta cómo
ese muro se resquebrajó en tres actos. Primero la grieta
(2024--2025): un distinguidor subexponencial y su réplica formal.
Después la primavera de 2026: ataques algebraicos que rompen retos
públicos y acercan las estimaciones, aún por encima de los niveles
declarados. Y por fin el verano: el método \emph{holdout}, un
distinguidor cuasipolinomial de origen insospechado, y una cascada de
estimaciones de recuperación de clave que caen a $2^{89}$--$2^{149}$
\emph{por debajo} de los niveles declarados de todos los conjuntos
---con supuestos heurísticos y sin ataque práctico alguno---. Cerraremos
con el debate técnico (¿cuánto cuesta mover los bits?), la respuesta
institucional (el equipo, el BSI, ISO, NIST) y una lectura sobria de
lo que un lector de este libro debe llevarse. La regla del capítulo:
separar a cada paso lo demostrado de lo heurístico, y lo heurístico de
lo titulado.
\end{roadmap}

\section{La primera grieta (2024--2025)}\label{sec:2026-grieta}

La tradición del capítulo~\ref{ch:estructura} tenía una moraleja
estadística: cada intento estructural se quedaba corto por miles de
bits. Esa estadística empezó a romperse en 2024, cuando Hugues
Randriambololona publicó el \emph{distinguidor de sizigias}
(«syzygy distinguisher»)~\cite{randriam2024}, publicado después en
EUROCRYPT~2025. Su hallazgo: las claves públicas de Goppa binaria se
pueden distinguir de matrices aleatorias con complejidad
\emph{subexponencial} en la capacidad de corrección del código. Para
\texttt{mceliece348864}, la estimación rondaba los $2^{529}$; para
los conjuntos recomendados por el BSI, al menos $2^{683}$ (véase su
evaluación de 2026~\cite{bsi2026}). Frente a los $2^{9000}$ de la
tabla~\ref{tab:est-espacios}, la grieta era visible; frente a los
$2^{140{,}8}$ del mejor ataque genérico, seguía siendo abisal.

La respuesta del equipo llegó el 17 de abril de 2025, en una nota
fechada titulada precisamente \emph{Notes on a recent claim that a
\texttt{mceliece348864} distinguisher uses only $2^{529}$
operations}~\cite{cm529}. Es una pieza ejemplar de metodología, y
conviene resumir sus tres argumentos porque enmarcan todo lo que
viene después.

\begin{enumerate}[leftmargin=1.6em,itemsep=3pt]
  \item \textbf{La seguridad no descansa en la indistinguibilidad.}
        La cadena del capítulo~\ref{ch:seguridad} va de arriba
        abajo: IND-CCA2 en el QROM se apoya en OW-CPA del cifrado,
        y OW-CPA en la dureza de la decodificación. Que la clave
        pública sea reconocible como Goppa no toca ningún eslabón:
        el equipo lo ilustra con RSA, cuyas claves se distinguen de
        enteros aleatorios de forma trivial sin que ello amenace la
        factorización~\cite{cm529}.
  \item \textbf{Hay un suelo trivial de $2^{256}$.} La clave privada
        se genera a partir de una semilla de 256 bits; un atacante
        puede, en el peor caso, enumerar semillas. Cualquier ataque
        de recuperación de clave por debajo de $2^{256}$ sería una
        noticia de primera magnitud; los $2^{529}$ no lo son. Esta
        observación, casi trivial, se convertirá en la vara de medir
        del capítulo.
  \item \textbf{El historial manda prudentemente.} La nota repasa la
        larga lista de «avances» algebraicos previos ---desde los
        ataques de 2002 que se anunciaban como amenaza para AES
        hasta un trabajo de 2010 que se describía como
        «breakthrough»--- todos ellos digeridos sin consecuencias:
        el método muere al comparar constantes, o no escala, o
        confunde distinguir con romper.
\end{enumerate}

Añadamos dos piezas de contexto. En marzo de 2025, NIST cerró su
cuarta ronda eligiendo HQC como respaldo basado en códigos y dejó a
Classic McEliece fuera de la estandarización, anotando en el
informe IR~8545 que la criptoanálisis reciente «debilita en cierta
medida el argumento de tratarlo como una opción especialmente
conservadora» y señalando que \texttt{mceliece460896} ya quedaba
corto frente a su categoría~\cite{nist8545}. Y en noviembre de 2025,
el equipo publicó una nota sobre \emph{You Only Decapsulate Once},
un ataque de canal lateral que explota variaciones de tiempo
introducidas por compiladores modernos en software criptográfico:
una advertencia de que la superficie de ataque de una implementación
no termina en su algoritmo~\cite{cm2025yodo}.

\begin{historia}[De 2002 a 2024: el largo cementerio algebraico]
La nota del equipo de abril de 2025~\cite{cm529} y la de junio de
2026~\cite{cm610} enumeran, con paciencia de archivero, las
promesas algebraicas caídas. La de 2026 abre con un ejemplo de 2002
que anunciaba coste «subexponencial» contra cifrados de bloque y
resultó inofensivo; sigue con la lista de trabajos contra McEliece
---2010: «breakthrough in code-based cryptography»; 2014:
reducción de Wild--Goppa; 2016: el método de claves compactas de
Faugère--Otmani--Perret--de~Portzamparc--Tillich; 2023: el
modelo Pfaffiano de Couvreur--Mora--Tillich; 2024: las sizigias de
Randriambololona---. La lección del cementerio no es que los papers
fueran malos: es que durante dos décadas cada funeral corporativo
se celebraba antes de la autopsia. En 2026, por primera vez, el
ritual se invertirá: primero las réplicas, después los títulos.
\end{historia}

\section{La primavera de 2026: de CFS a los retos TII}\label{sec:2026-primavera}

El 3 de marzo de 2026 apareció el primer golpe con \emph{víctima
real} fuera de los parámetros de Classic McEliece: Bardet, Lemoine y
Tillich rompieron el esquema de firmas CFS, de 25 años de
antigüedad, construido sobre códigos de Goppa binarios de tasa muy
alta, aprovechando el modelo Pfaffiano de 2023 y una búsqueda de
cuadrados en el ideal correspondiente. El mismo documento recuperó
las claves de retos de TII «con seguridad declarada de hasta 210
bits»~\cite{bardet2026}. La letra pequeña salvó a Classic McEliece:
el ataque vive de tasas de código extremas, y los conjuntos del
esquema tienen tasas de $0{,}7$--$0{,}8$. Pero el cementerio acababa
de enterrar un esquema entero.

En junio llegaron dos piezas más. El 10 de junio, Briaud, Lemoine,
Randriambololona y Tillich (BLRT) publicaron un ataque algebraico
\emph{subexponencial heurístico} basado en encontrar matrices de
rango~2 en el código de relaciones cuadráticas asociado a la Goppa;
con ellas se recupera la estructura secreta y, de ahí, una clave
equivalente. Rompieron retos TII de 83, 89, 119, 166, 210 y 248
bits declarados, y una clave CFS en 14 horas y 24\,GB de
RAM~\cite{briaud2026}. Para los parámetros de Classic McEliece
nivel~1, los propios autores cifran el ataque en «a lo sumo del
orden de $2^{454}$». El 23 de junio, el equipo respondió con una
nota fechada~\cite{cm610}: auditando el cómputo, el coste varía con
las suposiciones algebraicas entre $2^{454}$ (Wiedemann), $2^{513}$,
$2^{515}$ (con exponente de multiplicación $2{,}373$), $2^{607}$ y
$2^{610}$; ninguna variante baja del suelo trivial $2^{256}$
(ninguno de esos trabajos lo compara, observa el equipo con
retintín); los retos TII rotos son juguetes ($n\le 1409$) que
«omiten una de las defensas centrales» al reducir $t$ y con él las
opciones de $g$; y la narrativa del «ataque que rompe el esquema»
repite, palabra por palabra, la de 1991.

El 29 de junio, Tobias Hemmert ---del propio BSI--- convirtió en
recuperación de clave el distinguidor de «anulación de orden
superior» (\emph{higher-order vanishing}) que él y Andreas Wiemers
habían presentado en CRYPTO~2026~\cite{hemmert2026,hemmertwiemers2026}:
los polinomios homogéneos que se anulan con multiplicidad alta en
las columnas de la matriz de paridad no solo distinguen Goppa de
aleatorio; su estructura permite reconstruir la clave. Conjetura
del autor: el coste sería «aproximadamente el mismo que el del
distinguidor», que el BSI evalúa en al menos $2^{550}$ para los
conjuntos que recomienda~\cite{bsi2026}. Markku-Juhani Saarinen
extendió el método (HOVER) sin mejora significativa para estos
parámetros~\cite{saarinen2026hover}. Balance de la primavera: las
estimaciones estructurales habían bajado de miles de bits a
$2^{454}$--$2^{742}$ ---inmensas, pero ya del orden de la decena
de miles de millones de veces «solo» por encima de los objetivos,
y con la dirección de la flecha en contra.

\begin{historia}[1991, 1997, 2026: la misma obra, el mismo papel]
La nota de junio de 2026~\cite{cm610} recuerda el precedente: un
«criptoanálisis del criptosistema de McEliece» de 1991, que resultó
ser un ataque sobre un juguete con $n=63$; la historia la contó
Berson en 1997, y este libro la citó en el capítulo~\ref{ch:kem}
como anécdota de época. El patrón se repite: una demostración en
tamaño de juguete, una extrapolación en el titular, una espera de
quince años hasta saber si el ataque es un embrión o un espejismo.
La diferencia de 2026 es la velocidad: lo que en 1991 tardó un año
en replicarse, aquí tarda semanas ---y la propia respuesta del
equipo está fechada con precisión notarial.
\end{historia}

\section{El método \emph{holdout} (agosto de 2026)}\label{sec:2026-holdout}

El 7 de agosto de 2026, Ashrujit Ghoshal, Yuval Ishai, Aayush Jain y
Nuozhou Sun (GIJS) publicaron un distinguidor \emph{cuasipolinomial}
para Goppa--McEliece en el régimen asintótico del esquema
($m=\Theta(\log n)$, $t=\Theta(n/\log n)$), con tiempo
$n^{\mathcal{O}(\log n)}$ y ventaja $1-o(1)$. El origen del método
es delicioso: nació de un intento \emph{fallido} de construir
protocolos de recuperación privada de información (PIR) a partir de
códigos localmente decodificables. El título del
artículo~\cite{ghoshal2026} lo confiesa: «PIR Meets McEliece».

La idea, en la figura~\ref{fig:2026-holdout}. Las columnas de la
clave pública son evaluaciones de polinomios secretos: yacen
---salvo reescalado--- sobre una curva algebraica oculta. El test
\emph{retira} (hold out) una columna y busca polinomios de bajo
grado que se anulen con multiplicidad alta en todas las demás. En
una clave real, esos polinomios también se anulan en la columna
retirada: es la estructura de curva forzándolo. En una matriz
aleatoria, solo ocurriría por azar. Ese sesgo, amplificado sobre
muchas retiradas, es el distinguidor. Así descrito parece poco,
pero trae la propiedad que ningún distinguidor previo tenía: la
complejidad crece \emph{cuasipolinomialmente} en el régimen de los
parámetros, y el método aplica a los cinco conjuntos, con
estimaciones de $2^{114}$--$2^{124}$ operaciones de bits para el
test (frente a los $2^{151}$--$2^{287}$ de la decodificación
genérica), aunque con almacenamiento de exabytes ($2^{66}$--$2^{71}$
bits) que los autores no intentaron optimizar~\cite{ghoshal2026}.

\begin{figure}[htbp]
\centering
\begin{tikzpicture}[
    col/.style={draw=mcblue,fill=mcsky!35,line width=0.4pt},
    colx/.style={draw=mcorange,fill=mcorange2!70,line width=0.4pt},
    dot/.style={circle,fill=mcred,inner sep=1.1pt},
    lbl/.style={font=\scriptsize,align=center},
    lblt/.style={font=\tiny,color=mcgray,align=center}
  ]
  % columnas de la clave publica
  \foreach \i in {0,1,2,3,4,5,6} {
    \draw[col] (0.95*\i,0) rectangle +(0.7,1.6);
  }
  \draw[colx] (7.1,0) rectangle +(0.7,1.6);
  \node[lblt] at (2.85,-0.25) {columnas de $G$ (clave pública)};
  \node[lblt,mcorange] at (7.45,-0.25) {retirada};
  \draw[-{Stealth[length=2mm]},mcorange,thick]
    (9.0,2.75) -- (7.85,2.45);
  \node[lbl,mcorange,text width=3.4cm] at (11.0,2.95)
    {se retira una columna y se busca un polinomio $P$ que se anule con
     multiplicidad alta en las demás};
  % puntos: P se anula
  \foreach \i in {0,1,2,3,4,5,6} {
    \node[dot] at (0.95*\i+0.35,1.85) {};
  }
  \node[lbl,mcred] at (3.0,2.15) {$P$ se anula en todas ellas};
  % la retirada: signo de pregunta
  \node[font=\Large,color=mcorange] at (7.45,1.95) {?};
  % paneles de respuesta
  \node[draw=mcgreen,fill=mcgreen2!45,rounded corners=3pt,lbl,
        text width=3.4cm,inner sep=5pt] at (11.9,1.15)
    {\textbf{clave Goppa real:}\\[2pt] $P$ también se anula en la
     retirada (lo fuerza la curva oculta)};
  \node[draw=mcred,fill=mcred2!45,rounded corners=3pt,lbl,
        text width=3.4cm,inner sep=5pt] at (11.9,-0.7)
    {\textbf{matriz aleatoria:}\\[2pt] se anula solo por azar; el
     test lo delata};
\end{tikzpicture}
\caption{El test \emph{holdout} de GIJS~\cite{ghoshal2026}: retirar
una columna y buscar polinomios de bajo grado que se anulen con
multiplicidad alta en todas las demás. Sobre una clave real el
polinomio se anula también en la retirada (estructura de curva);
sobre una matriz aleatoria, no. La extensión a recuperación de
clave usa la geometría local de la curva en la columna
retirada~\cite{vedenev2026,weis2026}.}
\label{fig:2026-holdout}
\end{figure}

El 27 de agosto, GIJS añadieron a su artículo los detalles de un
ataque de \emph{recuperación de clave} heurístico basado en la
misma maquinaria, con estimaciones de $2^{130{,}4}$ a $2^{149{,}0}$
según el conjunto, y de un ataque de recuperación de texto plano
que el equipo evaluó en unos $2^{266}$ para
\texttt{mceliece6960119}~\cite{ghoshal2026,cm610}. La revisión
también declaraba, con honestidad ejemplar, que el ataque «no es
todavía concretamente eficiente» pero que «puede ser relevante para
los niveles de seguridad de NIST». En paralelo, la réplica pública
del equipo (21 de agosto) recordaba que un distinguidor no ataca
los objetivos de seguridad del esquema ---el argumento de
abril de 2025--- y que la maquinaria no produce, por ahora, ningún
ataque ejecutable~\cite{bernstein2026,cm529}.

\section{La cascada de septiembre}\label{sec:2026-cascada}

Lo que ocurrió después fue una carrera de tres semanas entre grupos
---con préstamos cruzados, porque todos parten del test de GIJS---
para llevar el \emph{holdout} de distinguidor a recuperación de
clave. Conviene registrar los hitos con fechas, porque la velocidad
es parte de la historia.

\begin{itemize}[leftmargin=1.5em,itemsep=3pt]
  \item \textbf{Vedenev (20 de agosto; revisión de octubre).} Extiende
        el método a la recuperación de clave para subcódigos de
        GRS (y por tanto alternantes y Goppa): un \emph{único} tipo
        de ejecución del test recupera la clave, con coste estimado
        de $2^{89{,}9}$ a $2^{98{,}6}$ operaciones de bits para los
        parámetros de Classic McEliece. Demostración: claves de los
        retos TII-173, TII-249 y TII-253, «los dos primeros no
        resueltos hasta donde sabemos»~\cite{vedenev2026}.
  \item \textbf{Saarinen (24 de agosto; revisiones).} Publica una
        contabilidad de costes del método («Holdout»), de $2^{126{,}8}$
        a $2^{145{,}2}$ según la revisión de septiembre, y resuelve
        el reto TII-254 ---el más difícil resuelto hasta la fecha
        en su métrica original, $2^{254}$--- con una variante de dos
        anclas: dos núcleos de relaciones de dimensión 121, un
        núcleo común de 80, una geometría proyectiva sobre
        $\GF{2^8}$, y una secuencia de Krylov final de 27{,}2
        horas-GPU en aceleradoras NVIDIA~\cite{saarinen2026a,saarinen2026b}.
  \item \textbf{Weis (11 de septiembre; versión 0.5 de octubre).}
        Stephen Weis, de Anthropic, mejora el ataque en dos frentes.
        Primero, el coste de una ejecución baja a $2^{89}$--$2^{98}$
        operaciones de bits (sin cargo de memoria); pagando cada
        operación por la memoria que direcciona, a coste de raíz
        cuadrada, $2^{107}$--$2^{117}$; bajo la regla contable de la
        guía de seguridad ---que carga el coste de toda la memoria
        en cada operación---, $2^{117}$--$2^{128}$, contra
        $2^{159}$--$2^{316}$ de la decodificación genérica en ese
        mismo modelo. Segundo, dos rutas de recuperación de clave:
        una apoyada en un teorema que demuestra y que exige entre
        100 y 1400 ejecuciones; otra, de una sola ejecución, que
        construye sobre el key recovery de GIJS. «Nada de esto está
        cerca de lo práctico ---escribe---, y una ejecución a tamaño
        completo no puede probarse»~\cite{weis2026}.
  \item \textbf{Apon (26--28 de agosto) y la réplica.} Daniel Apon
        publica una cota inferior de la ruta de Vedenev: por razones
        algebraico-geométricas, el número de columnas retiradas que
        su paso crítico necesita es mayor que $2t+3$; para
        \texttt{mceliece8192128}, eso implica un coste superior a
        $2^{1500}$ operaciones de bits. Vedenev respondió con una
        revisión que añade una segunda ruta
        («código menor») y sostiene que las objeciones ya no aplican
        a la versión nueva~\cite{apon2026,vedenev2026}.
\end{itemize}

La tabla~\ref{tab:2026-costes} reúne las estimaciones publicadas con
el mismo protocolo: operaciones de bits, sin cargo de memoria, para
los cinco conjuntos. Es la tabla que en 2022 habría sido ciencia
ficción: la columna estructural ya no está en los miles de bits.

\begin{table}[htbp]
\centering
\footnotesize
\setlength{\tabcolsep}{6pt}
\begin{tabular}{lccccc}
\toprule
conjunto & decodificación & GIJS & Saarinen & Weis &
AES ref.\\
 & (ISD, cap.~\ref{ch:parametros}) & (27-ago) & (15-sep) & (v0.5) & (NIST)\\
\midrule
\texttt{mceliece348864} & $2^{151{,}4}$ & $2^{130{,}4}$ & $2^{126{,}8}$ & $2^{94{,}3}$ & $2^{143}$\\
\texttt{mceliece460896} & $2^{193{,}3}$ & $2^{149{,}0}$ & $2^{145{,}2}$ & $2^{99{,}3}$ & $2^{207}$\\
\texttt{mceliece6688128} & $2^{268{,}0}$ & $2^{141{,}3}$ & $2^{137{,}5}$ & $2^{102{,}4}$ & $2^{272}$\\
\texttt{mceliece6960119} & $2^{268{,}4}$ & $2^{140{,}5}$ & $2^{136{,}7}$ & $2^{101{,}6}$ & $2^{272}$\\
\texttt{mceliece8192128} & $2^{303{,}2}$ & $2^{141{,}6}$ & $2^{137{,}5}$ & $2^{102{,}4}$ & $2^{272}$\\
\bottomrule
\end{tabular}
\caption{La fotografía de octubre de 2026, en operaciones de bits y
sin cargo de memoria. La columna ISD es el mejor ataque genérico sin
trucos de memoria (guía de seguridad; capítulos~\ref{ch:parametros}
y~\ref{ch:isd}). Las columnas de 2026 son estimaciones
\emph{condicionales}, apoyadas en supuestos heurísticos, de ataques
de recuperación de clave: GIJS (revisión del 27 de agosto, su
tabla~3), Saarinen (revisión del 15 de septiembre) y Weis (versión
0.5 de octubre, tabla~1)~\cite{ghoshal2026,saarinen2026a,weis2026}.
La última columna es la referencia de NIST para las categorías
1/3/5 (búsqueda de clave sobre el AES correspondiente). Ninguna de
las cifras de 2026 corresponde a un ataque ejecutado ni ejecutable
en hardware real.}
\label{tab:2026-costes}
\end{table}

Un matiz imprescindible, del propio equipo: los retos TII resueltos
---$n=214$, $n=223$--- tienen $t$ diminutos y por tanto pocas
opciones de polinomio $g$; «omiten una de las defensas centrales»
del esquema~\cite{cm610}. Que la maquinaria digiera juguetes no
prueba que digiera los conjuntos reales: eso lo dirá el tiempo, o
una ejecución a escala que nadie ha hecho.

\section{El debate de los costes}\label{sec:2026-costes}

Los números de la tabla~\ref{tab:2026-costes} comparten un supuesto
silencioso: que una operación de bits cuesta lo mismo que cualquier
otra. Daniel Bernstein lo atacó de frente en el foro de NIST. Un
\emph{run} del ataque es una cadena de $2^{32}$ a $2^{36}$ pasadas
mutuamente dependientes sobre un array de 28 a 540~TiB
($2^{47{,}8}$--$2^{52{,}1}$ bits); el ataque contra
\texttt{mceliece348864} movería del orden de $2^{93}$ bits entre
memoria y unidades aritméticas. Enrutar esos bits a la velocidad de
interconexión de los superordenadores actuales, calculaba Bernstein,
llevaría unos 90\,000 años, solo para esa componente; y las cotas
de área-tiempo de Brent--Kung sugieren que ningún trazado de
circuito elimina el peaje. La misma asimetría se esconde en la
comparación con AES: la búsqueda de claves no toca memoria, se
amortiza entre muchos objetivos y gana con Grover; el
\emph{holdout} no hace nada de eso~\cite{bernstein2026}.

Weis, lejos de esquivar el golpe, lo incorpora: su artículo admite
que no conoce distribución de datos que haga local el acceso, y
ofrece la escalera completa de modelos contables. Con cargo de
raíz cuadrada ---un modelo con precedente en la propia guía---, una
ejecución cuesta $2^{107}$--$2^{117}$: aún por debajo de la
decodificación genérica con memoria libre. Con la regla más severa
de la guía, $2^{117}$--$2^{128}$. La discusión no es bizantina: es
\emph{el mismo} debate metodológico del capítulo~\ref{ch:isd},
ahora aplicado a la otra orilla. Y tiene precedente ilustre: en
2020, el propio Bernstein argumentó que NIST aplicó a Kyber-512 la
vara de las compuertas clásicas ($2^{143}$), y NIST respondió que
tiende a dar el beneficio de la duda a un conjunto cuyos ataques
quedan a pocos bits de la categoría. Aquí la distancia no es de
pocos bits: bajo el modelo más severo, Weis estima 25 bits bajo la
categoría~1 para \texttt{mceliece348864} y 82--145 bits bajo las
categorías del resto de conjuntos; y Apon declaró el 1 de octubre
tener un análisis no publicado que sitúa el nivel~1 en
$2^{89}$--$2^{92}$ \emph{incluyendo} costes de memoria, pero
decidió no publicarlo «para no perjudicar a los usuarios
actuales»~\cite{weis2026,bernstein2026}.

\begin{rigor}[Demostrado, heurístico, estimado: el mapa de 2026]
La contabilidad honesta de esta historia tiene tres columnas.
(1)~\textbf{Demostrado}: el distinguidor de GIJS está probado (en
su régimen asintótico), y Weis prueba varios enunciados sobre
Goppa binaria que sostienen su ruta. (2)~\textbf{Heurístico}: la
recuperación de clave a tamaño real depende de supuestos de rango
y genericidad ---Weis enumera cinco, los prueba en claves pequeñas
y en retos TII ($n=214$ frente a $n\ge3488$) y argumenta por qué
las desviaciones observadas no ocurrirían a escala---; el grado
del ataque es 7; el paso caro (el cómputo de núcleos) no se ha
ejecutado contra ningún conjunto real. (3)~\textbf{Estimado}: todas
las cifras de la tabla~\ref{tab:2026-costes} dependen de un modelo
de costes ---el de las operaciones de bits desnudas--- que el
propio equipo disputa, y que mueve las cuentas 20--30 bits según
el cargo. Nada de la columna (2)~(3) convierte a 2026 en «McEliece
roto»; nada de la columna (1) permite ya decir «McEliece inmune».
El estado del arte es exactamente eso: un arte, con supuestos
declarados y réplicas en curso.
\end{rigor}

\section{La respuesta institucional}\label{sec:2026-institucional}

El equipo de Classic McEliece sostuvo durante todo el otoño una
posición pública en dos capas. Como equipo: la réplica del 21 de
agosto al distinguidor ---no ataca los objetivos, no hay ataque
ejecutable--- y la evaluación del ataque de texto plano
($\approx 2^{266}$ para \texttt{mceliece6960119}). Y la respuesta
fechada de junio~\cite{cm610}. Como individuos: Bernstein mantuvo,
el 30 de septiembre, que ninguna nación-estado tiene recursos para
ejecutar los ataques, que estos no ganan nada con múltiples
objetivos ni con computación cuántica, y que su evaluación relativa
---ML-KEM con más riesgo que McEliece--- no había cambiado. La
pregunta directa de Filippo Valsorda ---¿sigue el equipo
recomendando sus conjuntos estandarizados?--- no había recibido
respuesta etiquetada como del equipo a fecha de cierre de este
libro~\cite{bernstein2026}.

El movimiento decisivo fue del BSI, el 1 de octubre de 2026, en una
nota de cuatro páginas fechada el 29 de septiembre~\cite{bsi2026}.
Su evaluación, en orden: los distinguidores previos no amenazaban
los supuestos (el de sizigias, $\ge 2^{683}$; el de
Hemmert--Wiemers, $\ge 2^{550}$; BLRT, $\ge 2^{742}$); los métodos
de agosto --octubre sí merecen acción, con cifras de $2^{120}$--$2^{124}$
(distinguidor GIJS) y $2^{140}$--$2^{148}$ (recuperación GIJS) para
los conjuntos de su guía, y $2^{99}$--$2^{102}$ en la contabilidad
de Weis; los resultados «no permiten todavía un ataque práctico»,
pero la mejora estructural «es significativa» y se esperan más.
Conclusión operativa: Classic McEliece «no debería seguir usándose
en nuevos desarrollos»; para intercambio de claves, preferir
ML-KEM o FrodoKEM, y HQC cuando su estándar aparezca; los
despliegues híbridos existentes conservan al menos la seguridad de
la mitad clásica; y la recomendación formal se revisará en
TR-02102-1 a principios de 2027. El BSI añadió una lección de
migración: los resultados «subrayan la importancia de la
\emph{crypto-agility}», y ninguno de los ataques toca a los otros
esquemas recomendados ---ni a HQC, que no usa códigos de Goppa ni la estructura
oculta que estos ataques explotan~\cite{bsi2026}.

En paralelo, el ecosistema: NIST dejó la puerta abierta a
estandarizar McEliece más adelante, pero su IR~8545 ya había
elegido HQC y anotado los recortes de \texttt{mceliece460896}; ISO
publicó en junio de 2026 su enmienda con los cuatro conjuntos
grandes ---sin \texttt{mceliece348864}, el más expuesto a los
nuevos métodos--- y cada uno de esos cuatro tiene ahora
estimaciones por debajo de su nivel declarado; las agencias
europeas divergen (ANSSI ya lo situaba por debajo de ML-KEM,
FrodoKEM y HQC; el manual holandés lo tenía por conservador; la
NCSC sueca solo bendice estándares NIST); y los despliegues reales
---túneles VPN de Mullvad, Rosenpass, PQConnect---, todos en modo
híbrido, activaron sus protocolos de revisión~\cite{bernstein2026,bsi2026,nist8545}.

\begin{trampa}[Los dos titulares fáciles]
Desde agosto de 2026 circulan dos frases gemelas, y ambas son
falsas. «\emph{McEliece ha caído}»: nadie ha recuperado una clave
ni descifrado un texto de ningún conjunto real; los propios
autores de los ataques lo repiten; el ataque más caro de ejecutar
---el cómputo de núcleos--- no se ha corrido a escala; los retos
rotos son juguetes. «\emph{McEliece es inmune}»: las estimaciones
de recuperación de clave de tres grupos independientes están por
debajo de los niveles que el propio esquema declaró, el BSI ---que
lo recomendaba desde 2020--- ha pedido detener nuevos despliegues,
y la incógnita no es si los supuestos heurísticos fallarán, sino
cuándo y cuánto. La versión utilizable de la verdad tiene tres
cláusulas y ninguna es un titular: nadie ha roto nada, muchas
estimaciones condicionales han bajado de los niveles declarados, y
la comunidad tarda años en digerir cuánto pesa cada supuesto. Quien
lea una sola de las dos frases, que relea este recuadro.
\end{trampa}

\section{Lectura final}\label{sec:2026-lectura}

¿Qué debe llevarse el lector de este libro de todo el terremoto?
Cuatro cosas, ninguna grandilocuente.

\textbf{Primera: la criptografía finita se juega en los márgenes.}
Los conjuntos de Classic McEliece llevaban cuarenta y ocho años
fuera de rango de todos los ataques estructurales conocidos, y
estaban elegidos contra \emph{un modelo} de ataque (la
decodificación genérica) con márgenes que la propia guía declaraba
para el otro modelo en miles de bits. El método \emph{holdout}
apareció fuera de ese modelo, traído de la teoría de PIR, y la
lección no es «los márgenes eran falsos» sino la del
capítulo~\ref{ch:seguridad}: \emph{seguridad es la mejor estimación
disponible más los supuestos, y los supuestos envejecen}.

\textbf{Segunda: el eje manda.} Weis midió que el coste del
\emph{holdout} crece unos ocho bits por cada unidad del grado $m$
del cuerpo, y apenas depende de $n$: por eso sus conjuntos
recomendados (todos con $m=12$ o $13$) se comportan en bloque. Sus
ilustraciones ---$m=18$, $26$, $34$, con claves de $0{,}6$, $2{,}1$ y
$7{,}1$ megabytes--- y la instancia de Bernstein ($m=18$,
$n=6176$, $t=112$, clave de~1\,MB, texto cifrado de 252 bytes,
ejecución estimada en $2^{142}$) dibujan la dirección de una
futura re-parametrización: cuerpos más grandes, $n$ bien por debajo
de $2^m$, claves más grandes. Nadie ha propuesto tal conjunto para
estandarización, y Weis lo etiqueta «ilustraciones»~\cite{weis2026}.

\textbf{Tercera: la asimetría de las pruebas.} El distinguidor está
probado; la recuperación de clave, no. Los experimentos existen
solo en tamaño de juguete. Quien compare las cotas del capítulo
con las de este debe recordar la dirección del rigor: una
estimación heurística baja no convierte a un esquema en roto ---ni
una superioridad histórica lo convierte en intacto. La única
moneda válida es la misma que usa este libro desde la
página~1: enunciar el modelo, citar la fuente, y separar lo
demostrado de lo conjeturado.

\textbf{Cuarta: esto es historia viva.} A la fecha de este libro
(octubre de 2026), las réplicas siguen circulando (la última
revisión de Vedenev es del 7 de octubre), el BSI revisará su guía a
principios de 2027, ISO mantiene publicado su estándar, y las
estimaciones de los distintos grupos ---dentro del mismo marco---
difieren en decenas de bits. Cualquier lector que llegue a estas
páginas años después debería buscar la actualización de la
guía~\cite{cmsecurity}, las respuestas fechadas del equipo
---la serie de notas continuará--- y las réplicas de las revistas
que aún están por llegar. El resto del libro no depende de la
resolución del terremoto: el mecanismo, los conjuntos, la ISO, el
laboratorio, todo lo anterior sigue siendo exactamente lo que las
páginas describen. Lo que cambia es la conversación sobre si
conviene usarlo, y esa conversación, por primera vez en la historia
del esquema, está abierta.

\section{Ejercicios}

\begin{enumerate}[leftmargin=1.6em,itemsep=4pt]
  \item \textbf{El muro de 2022 contra el suelo de 2026.}
        (a)~Compare la mejor estimación estructural de la
        tabla~\ref{tab:est-espacios} para \texttt{348864}
        ($2^{9116{,}8}$) con la de Weis ($2^{94{,}3}$): calcule la
        diferencia de exponentes y comente en dos líneas por qué
        \emph{no} es el «mismo ataque mejorado» (pista: ¿qué
        familia de métodos estaba agotada y cuál es nueva?).
        (b)~Compare ahora $2^{94{,}3}$ con el suelo trivial $2^{256}$
        de la respuesta del equipo~\cite{cm529}: ¿qué implica que
        la estimación esté por debajo?
  \item \textbf{Una ejecución, cien ejecuciones.} Con la ruta de
        Weis que exige entre 100 y 1400 ejecuciones, calcule el
        coste total si una ejecución cuesta $2^{89}$ (sin memoria)
        y si cuesta $2^{117}$ (con cargo severo). Encaje ambos
        resultados con la afirmación «esta es nuestra mejor
        estimación; con el cargo de la guía queda por debajo de la
        decodificación genérica en ese modelo».
  \item \textbf{El eje m.} Partiendo de $2^{94{,}3}$ ($m=12$) y de
        la regla de Weis (unos 8 bits por unidad de $m$), estime el
        coste para $m=18$ y compárelo con la instancia de Bernstein
        ($n=6176$, $t=112$), cifrada en $\approx 2^{142}$.
        ¿Cuánto sube la clave pública (en MB) frente a las actuales?
  \item \textbf{Distinguir, recuperar, romper.} Clasifique cada
        afirmación como demostrada, heurística o estimada, usando
        el recuadro de rigor: (a)~«las claves de Goppa binaria se
        distinguen de matrices aleatorias con ventaja
        $1-o(1)$»; (b)~«recuperar la clave de \texttt{8192128}
        cuesta $2^{102{,}4}$»; (c)~«el ataque de Vedenev cuesta
        esencialmente una ejecución»; (d)~«nadie ha ejecutado el
        cómputo de núcleos contra un conjunto real».
  \item \textbf{El debate de la memoria, aplicado.} Use la
        analogía de AES del capítulo~\ref{ch:isd} (tabla de
        equivalencias) y explique, en cinco líneas, por qué
        comparar $2^{94{,}3}$ con $2^{143}$ (AES-128) favorece al
        ataque, mientras que comparar $2^{117}$--$2^{128}$ con
        $2^{143}$ favorece al esquema. ¿Qué información falta para
        cerrar la comparación en un sentido o en otro?
  \item \textbf{Los titulares.} Escriba su propio titular de
        longitud máquina (menos de 80 caracteres) para la noticia
        del 1 de octubre de 2026, que sea \emph{verdadero} según el
        recuadro de trampa, y otro deliberadamente falso en cada
        dirección. Intercámbielos con un colega y discutan qué
        cláusula falta en cada uno.
\end{enumerate}

\begin{center}\rule{0.55\textwidth}{0.4pt}\end{center}
\medskip
\noindent\textit{Cierre de la parte~\ref{part:seguridad}.} Cuatro
capítulos para responder a la pregunta «¿qué significa romper
McEliece?»: los objetivos y la cadena de reducciones
(capítulo~\ref{ch:seguridad}); la trinchera de la decodificación
genérica y sus pliegos de costes (capítulo~\ref{ch:isd}); el muro
estructural y su anatomía (capítulo~\ref{ch:estructura}); y el
terremoto de 2025--2026 con su balance aún abierto
(capítulo~\ref{ch:2026}). Nadie ha roto Classic McEliece ---no
existe un ataque práctico, y cada estimación baja de los niveles
declarados viaja con supuestos explícitos--- pero las cifras de esa
orilla han cambiado de siglo en cuatro años, y con ellas la
conversación pública. La parte final del libro vuelve al suelo: las
implementaciones, el ecosistema que ya existe, y el balance de lo
que este esquema significa para quien lee.
